\section{Introduction}

CIFAR-10 is a small-data benchmark that modern architectures have largely
saturated: with tens of millions of parameters and long schedules, test accuracy
above 99\% is routine. The problem changes when the model must be trained from
scratch under a hard budget of five million trainable parameters, and when the
evaluation set is hidden and may not match the training distribution. Both
constraints were active in this competition. The first forces careful
architecture and optimization choices; the second makes validation accuracy an
imperfect guide and turns the leaderboard into a noisy, partially observed
signal.

We approached the problem as a two-track collaboration. One track pushed clean
accuracy through architecture search, long schedules, and knowledge distillation
(KD) from large teachers trained on the same data. The other track studied
robustness: it characterized the only unlabeled images available to us, built
corrupted and photometric proxy validation sets from the training split, and ran
controlled augmentation ablations. The two tracks shared one codebase, one fixed
5,000-image holdout, and one set of conventions, so every number in this paper is
directly comparable.

The main result is a 4.88M-parameter student that reaches 97.84\% clean
validation accuracy and 97.92\% with 10-view test-time augmentation (TTA). The team's selected
submissions scored 0.911 and 0.904 on the public leaderboard; the team's private
score was 0.880, placing 13th. The roughly ten-point gap between
our best holdout number and the private score is the central unsolved problem of
the project.

An important caveat shapes how we interpret that gap. The 2,000 unlabeled images
distributed with the starter data, which we analyzed at length, were not the
images used for the final leaderboard. Our photometric statistics, nearest-
neighbor tests, and visual comparisons therefore describe that released sample,
not the final test set. They motivated the robustness experiments and gave us a
useful analysis pipeline, but they cannot explain the final score. We report them
as exploratory evidence and explicitly separate them from the leaderboard
outcome. The only direct evidence of a distribution difference is the gap between
validation and leaderboard accuracy itself.

This paper makes four contributions.
\begin{itemize}
\item A strong, reproducible 4.88M-parameter recipe: long cosine training, SAM,
EMA, TrivialAugment, Cutout, mixup/CutMix, and three-teacher KD. It improves a
plain baseline from 94.76\% to 97.84\% validation (97.92\% with TTA) and yields
the team's 0.911 public / 0.880 private submissions.
\item A controlled account of which ingredients mattered. Long schedules and KD
gave the largest gains; architecture changes within the 5M budget gave small ones;
several plausible tricks did not pay off.
\item A clean-neutral augmentation result: AugMix leaves clean accuracy unchanged
but improves a composite corruption proxy by 2.2pp and a photographic-jitter proxy
by 1.9pp, replicated across seeds and surviving distillation.
\item A negative-results and uncertainty report. We document failed techniques
with paired statistics and validation curves, and we separate proxy evidence from
evidence about the hidden test set.
\end{itemize}

Section~\ref{sec:related} reviews related work. Section~\ref{sec:setup} defines
the protocol. Section~\ref{sec:method} describes the architectures, training
recipe, KD variants, augmentation policies, and proxy analyses.
Section~\ref{sec:results} presents accuracy and robustness results.
Section~\ref{sec:negative} collects the techniques that underperformed.
Section~\ref{sec:discussion} discusses the gap and limitations;
Section~\ref{sec:conclusion} concludes.
